\section{Uniform local flatness}\label{sec:flatness}

This section proves the local estimates used to obtain uniform affine approximation in Section~\ref{sec:reduction}. The decisive estimate controls midpoint defects at every separation, including separations much smaller than the scale on which a function is known to be close to an affine function.

\subsection{A linear lift and affine blow-ups}

\begin{lemma}\label{lem:linear-lift}
If $u$ is infinity-harmonic in $\Omega\subseteq\R^d$, then
\[
  U(y,b)=u(y)+b
\]
is infinity-harmonic in $\Omega\times\R$.
\end{lemma}

\begin{proof}
Suppose that $\phi$ touches $U$ from above at $(y_0,b_0)$. Restriction to the vertical line gives $\phi_b(y_0,b_0)=1$. Put $p=\nabla_y\phi(y_0,b_0)$. If $p\ne0$, choose an affine function $\ell$ with $\ell(y_0)=b_0$ and $\nabla\ell\cdot p=1$. The function
\[
  \psi(y)=\phi(y,\ell(y))-\ell(y)
\]
touches $u$ from above at $y_0$. At that point $\nabla\psi=p$, because $\phi_b=1$, and
\[
  D^2\psi[p]
  =\phi_{yy}[p]+2\phi_{yb}\cdot p+\phi_{bb}
  =D^2\phi[(p,1)]=\Dinf\phi.
\]
The viscosity inequality for $u$ gives the required sign. If $p=0$, the vertical upper contact gives $\phi_{bb}\ge0$, and this is $\Dinf\phi$ because $\nabla\phi=(0,1)$. Lower tests are identical with the inequalities reversed.
\end{proof}

We recall the cone comparison and slope estimates of Crandall, Evans, and Gariepy \cite{CrandallEvansGariepy2001}. These give local Lipschitz regularity and the information about blow-ups used here. For an infinity-harmonic function $w$ on an open set $D\subseteq\R^n$, define, whenever $\overline B_r(z)\subset D$,
\[
  L_r^\pm(w;z)=\frac1r\max_{|h|\le r}\bigl\{\pm(w(z+h)-w(z))\bigr\}.
\]

\begin{lemma}\label{lem:cone-slopes}
The function $w$ is locally Lipschitz. For a fixed center $z$, both $L_r^+(w;z)$ and $L_r^-(w;z)$ are nondecreasing in $r$. For each fixed allowable radius, they are continuous as functions of the center.
\end{lemma}

\begin{proof}
Fix $\overline B_r(z)\subset D$ and $L>L_r^+(w;z)$. Choose $\gamma>0$ so small that
\[
  \gamma r<L-L_r^+(w;z),\qquad 2\gamma r<L.
\]
On $\partial B_r(z)$ the function $w$ lies strictly below
\[
  w(z)+L|Z-z|-\gamma|Z-z|^2.
\]
It lies below this function throughout the ball. Otherwise the positive maximum of their difference is attained at an interior point different from $z$, and the displayed function, shifted by a constant, is an upper test there. Writing $\rho=|Z-z|$, its infinity Laplacian is
\[
  -2\gamma(L-2\gamma\rho)^2<0,
\]
contrary to the viscosity inequality. Letting $\gamma\downarrow0$ and then $L\downarrow L_r^+(w;z)$ gives
\[
  w(Z)-w(z)\le L_r^+(w;z)|Z-z|\qquad (Z\in\overline B_r(z)).
\]
Applying the same argument to $-w$ gives the lower bound with $L_r^-(w;z)$. These bounds prove monotonicity in the radius.

For centers in a compact neighborhood, choose one small radius whose balls remain in $D$. The two slopes at that radius are uniformly bounded by the oscillation on a larger compact set divided by the radius. The preceding inequalities then give a common local Lipschitz bound. Continuity in the center at fixed radius follows from uniform continuity of $w$ on a compact neighborhood, taking the maximum over the fixed set $\{|h|\le r\}$.
\end{proof}

For the following fixed-center affine blow-up result, see Crandall, Evans, and Gariepy \cite{CrandallEvansGariepy2001} and Crandall and Evans \cite[pp.~124--127]{CrandallEvans2001}; we include a variant of the latter cone-slope argument.

\begin{proposition}[Affine blow-ups at a fixed center]\label{prop:affine-blowups}
Let $w$ be infinity-harmonic in $D$ and let $z\in D$. Every locally uniform limit, along a sequence $r_i\downarrow0$, of
\[
  V_i(Y)=\frac{w(z+r_iY)-w(z)}{r_i}
\]
is a linear function on $\R^n$. Such limits exist after passage to subsequences.
\end{proposition}

\begin{proof}
The local Lipschitz bound and the exhaustion of $\R^n$ by the rescaled domains give subsequential convergence on compact sets. By Lemma~\ref{lem:cone-slopes}, the finite limits
\[
  L_*^\pm=\lim_{r\downarrow0}L_r^\pm(w;z)
\]
exist. They are equal. Indeed, if $h_r$ attains the positive maximum, then for each fixed small $R>r$, once $r$ is sufficiently small,
\[
  L_r^+(w;z)\le L_r^-(w;z+h_r)\le L_R^-(w;z+h_r).
\]
First let $r\downarrow0$, using continuity at fixed $R$ and $|h_r|\le r$, and then let $R\downarrow0$. This gives $L_*^+\le L_*^-$. Reversing signs gives equality; write the common value as $L_*$.

Let $V$ be a limit. Uniform convergence on balls shows that
\[
  L_R^+(V;0)=L_R^-(V;0)=L_*\qquad(R>0),
\]
and that the slopes at every other center retain monotonicity in their radius. In particular $|V(Y)|\le L_*|Y|$. For any center $q$,
\[
  L_T^\pm(V;q)\le\frac{L_*(T+|q|)+|V(q)|}{T}.
\]
For a fixed radius at $q$, bound its slope by the slope at $T$ and send $T\to\infty$. Thus every slope of $V$ is at most $L_*$, so $V$ is globally $L_*$-Lipschitz.

If $L_*=0$, then $V=0$. Otherwise let $x_R$ and $y_S$ attain the positive and negative extrema on the balls of radii $R$ and $S$ centered at zero. The equality of the origin slopes gives
\[
  L_*(R+S)=V(x_R)-V(y_S)
     \le L_*|x_R-y_S|\le L_*(R+S).
\]
Equality throughout forces $|x_R|=R$, $|y_S|=S$, and these points to lie on opposite rays of one line. Comparing all choices of $R,S$ gives a unit vector $e$ such that $V(te)=L_*t$ for every real $t$. The Lipschitz bound then yields
\[
  L_*t-L_*|Y-te|\le V(Y)\le L_*t+L_*|Y-te|.
\]
Letting $t\to+\infty$ in the lower bound and $t\to-\infty$ in the upper bound gives $V(Y)=L_*e\cdot Y$.
\end{proof}

This argument establishes the linearity of each blow-up without requiring uniqueness of its slope. The stronger everywhere differentiability theorem is due to Evans and Smart \cite{EvansSmart2011}.

\subsection{The three-point estimate}

We use one standard viscosity tool. For a continuous function $g$, its second-order superjet $J^{2,+}g(z)$ consists of pairs $(p,X)$ for which
\[
  g(z+h)\le g(z)+p\cdot h+\tfrac12X[h]+o(|h|^2).
\]
The closed superjet $\overline J^{2,+}g(z)$ consists of limits of such jets at points approaching $z$, with the function values also converging. For the continuous operators used here, the viscosity inequality at upper tests applies to closed superjets as well.

\begin{lemma}[Theorem of sums, {\cite[Theorem~3.2]{CIL1992}}]\label{lem:sum-rule}
Let $g_j$ be continuous on open subsets of Euclidean spaces. Suppose that $\sum_{j=1}^N g_j(z_j)-\Psi(z_1,\ldots,z_N)$ has a local maximum at an interior product point, where $\Psi$ is $C^2$ near that point. Set $A=D^2\Psi$. For every $\eta>0$ there are symmetric matrices $X_j$ such that
\[
  (D_j\Psi,X_j)\in\overline J^{2,+}g_j(z_j),\qquad
  \operatorname{diag}(X_1,\ldots,X_N)\le A+\eta A^2.
\]
\end{lemma}

\begin{proposition}[Three-point flatness]\label{prop:three-point}
Let $U_i$ be infinity-harmonic on $B_1\subset\R^n$ with a common Lipschitz bound, and let $\ell_i$ be affine functions such that
\[
  \|U_i-\ell_i\|_{L^\infty(B_1)}\longrightarrow0.
\]
Then
\begin{equation}\label{eq:three-point}
  \sup_{\substack{x,y\in B_{1/4}\\x\ne y}}
  \frac{|U_i(x)+U_i(y)-2U_i((x+y)/2)|}{|x-y|}
  \longrightarrow0.
\end{equation}
\end{proposition}

\begin{proof}
Put $w_i=U_i-\ell_i$ and $\varepsilon_i=\|w_i\|_{L^\infty(B_1)}$. The slopes of $\ell_i$ are bounded: evaluating at $e/2$ and $-e/2$ for every unit vector $e$ gives $|\nabla\ell_i\cdot e|\le\Lip(U_i)+2\varepsilon_i$. Thus $w_i$ have a common Lipschitz bound $C$.

Suppose that \eqref{eq:three-point} fails. Passing to a subsequence and replacing $U_i,\ell_i$ by their negatives if necessary, there are endpoint pairs in $B_{1/4}$ with positive midpoint defect greater than $\delta|x-y|$ for some fixed $\delta>0$. For triples $x,y,z\in\overline B_{3/4}$ with $x\ne y$, write
\[
  c=\frac{x+y}{2},\qquad h=z-c,\qquad r=|x-y|,
  \qquad f(r)=r-\frac18r^{3/2}.
\]
Here $0<r\le3/2$, so $\nu r\le f(r)\le r$ with $\nu=1-\sqrt{3/2}/8>0$. Choose a smooth function $B(c)$ that equals $\delta/2$ on $\overline B_{1/4}$, is at least $\delta/2$ everywhere, and equals a sufficiently large constant $B_{\mathrm{out}}$ when $|c|\ge1/2$. Consider
\begin{equation}\label{eq:three-point-maximum}
  w_i(x)+w_i(y)-2w_i(z)-B(c)f(r)-\frac{|h|^2}{r^2}.
\end{equation}
Allowing the third point to move permits the use of Lemma~\ref{lem:sum-rule}, while the quadratic penalty keeps it close to the midpoint. The factor $B(c)$ localizes the maximum, and the concavity of $f$ will supply the dominant negative second variation.
The violating pair with $z=c$ makes this expression positive. On the other hand it is at most
\[
  Cr+2C|h|-B(c)f(r)-\frac{|h|^2}{r^2}
  \le Cr+C^2r^2-B(c)f(r).
\]
Its upper limit as $r\downarrow0$ is therefore at most zero, uniformly in the other variables. A positive maximum is attained at some triple with $r>0$.

At these maximizing triples, $r\to0$ as $i\to\infty$: positivity and $Bf\ge\delta\nu r/2$ give $r\le8\varepsilon_i/(\delta\nu)$. Comparing the maximizing value with the admissible choice $z=c$ also gives
\[
  \frac{|h|^2}{r^2}\le2\bigl(w_i(c)-w_i(z)\bigr)\le2C|h|,
  \qquad |h|\le2Cr^2.
\]
Choose $B_{\mathrm{out}}>(C+\tfrac32C^2)/\nu$. The preceding upper bound for \eqref{eq:three-point-maximum} is then nonpositive when $|c|\ge1/2$. Hence $|c|<1/2$ at a positive maximum. Since $r\to0$ and $h=O(r^2)$, the maximizing triples are interior to $(B_{3/4})^3$ for all large $i$.

At one such triple set $q=\nabla\ell_i$ and $e=(x-y)/r$. Apply Lemma~\ref{lem:sum-rule} to the functions $U_i,U_i,-2U_i$ and the test
\[
  \Psi=\ell_i(x)+\ell_i(y)-2\ell_i(z)+B(c)f(r)+\frac{|h|^2}{r^2}.
\]
This test is $C^2$ near the triple. Define $p_1=D_x\Psi$, $p_2=D_y\Psi$, and $p_0=-D_z\Psi/2$. Direct differentiation gives
\begin{align*}
  p_1&=q+\tfrac12f(r)DB(c)-h/r^2+a e,\\
  p_2&=q+\tfrac12f(r)DB(c)-h/r^2-a e,\\
  p_0&=q-h/r^2,\qquad
  a=B(c)f'(r)-2|h|^2/r^3.
\end{align*}
These vectors and $a$ are bounded. Since $f'(r)=1-3\sqrt r/16$ and $|h|=O(r^2)$, we also have $a\ge\delta/4$ for all sufficiently large $i$.

Let $\sigma=(p_1,p_2,p_0)$. The third jet supplied by Lemma~\ref{lem:sum-rule} has gradient $-2p_0$. The function $-2U_i$ is infinity-harmonic, so the viscosity inequalities for the three superjets give
\[
  X_1[p_1]\ge0,\qquad X_2[p_2]\ge0,\qquad X_3[-2p_0]\ge0.
\]
In particular $X_3[p_0]\ge0$. If $A=D^2\Psi$ at the maximum, the matrix inequality evaluated at $\sigma$ gives
\[
  0\le\operatorname{diag}(X_1,X_2,X_3)[\sigma]
       \le A[\sigma]+\eta|A\sigma|^2.
\]
The maximum, $A$, and $\sigma$ are fixed while $\eta\downarrow0$, although the jets may vary. It follows that $A[\sigma]\ge0$.

We compute the same quadratic form by moving $(x,y,z)$ along the straight line with velocity $\sigma$ and parameter $\tau$, with $\tau=0$ at the maximum. The endpoint relative velocity $p_1-p_2=2ae$ is parallel to $x-y$, and the third velocity differs from the mean endpoint velocity by $O(r)$. Write $c(\tau),h(\tau),r(\tau)$ for the induced quantities, and use dots for their derivatives at zero. The formulas above give
\[
  \dot r=2a,\qquad \ddot r=0,\qquad
  \dot c=\tfrac12(p_1+p_2)=O(1),\qquad
  \dot h=p_0-\dot c=-\tfrac12f(r)DB(c)=O(r).
\]
Since $h=O(r^2)$,
\[
  \left.\frac{d^2}{d\tau^2}\frac{|h(\tau)|^2}{r(\tau)^2}\right|_{\tau=0}
  =\frac{2|\dot h|^2}{r^2}
   -\frac{8\dot r\,h\cdot\dot h}{r^3}
   +\frac{6\dot r^{\,2}|h|^2}{r^4}=O(1).
\]
Also $f''(r)=-3/(32\sqrt r)$, and therefore
\begin{align*}
  \left.\frac{d^2}{d\tau^2}\bigl(B(c(\tau))f(r(\tau))\bigr)\right|_{\tau=0}
  &=B(c)f''(r)\dot r^{\,2}+2f'(r)\dot r\,DB(c)\cdot\dot c\\
  &\quad+f(r)D^2B(c)[\dot c]\\
  &=-\frac{3B(c)a^2}{8\sqrt r}+O(1)\longrightarrow-\infty.
\end{align*}
The affine part of $\Psi$ has zero Hessian. Thus $A[\sigma]\to-\infty$, contradicting its nonnegativity and proving \eqref{eq:three-point}.
\end{proof}

The three-point estimate is uniform over every nonzero endpoint separation.
This allows the observation scale and the scale of affine approximation
to tend to zero at different rates. Lemma~\ref{lem:family-flatness} uses
that freedom to obtain affine approximation uniformly over the full
normalized family of solutions and all interior centers.
